\chapter{NUMERICAL RESULTS}
\label{chap:results}

This chapter reads the experiments built in Chapter~\ref{chap:methodology} and lets them speak in a single arc. The arc begins with what each model costs to train (Section~\ref{sec:r_cost}), then pins down the one behaviour they all share by carrying spectral bias from an exact toy calculation to the kernel of a real network (Section~\ref{sec:r_bias}). With the mechanism established, the arc follows the error through the frequency bands as training proceeds (Section~\ref{sec:r_bands}), fixes a common parameter to compare the final spectra fairly (Section~\ref{sec:r_spectrum}), and isolates the one place where the physics term pays for itself (Section~\ref{sec:r_gibbs}). Two stress tests then push the operators off their training conditions (Section~\ref{sec:r_stress}), and Section~\ref{sec:r_synthesis} draws the threads together. The pseudospectral solver is the ground truth throughout, and every model leaves a failure mode that the arc makes explicit.

% ===============================================================
\section{The Price of Each Approach}
\label{sec:r_cost}

Before asking how well the models perform, we ask what they cost. Table~\ref{tab:training} lists the size, wall-clock training time and final training loss of the six trained models, and two axes organize it. Along the first, capacity, the operators carry about $2.1$ million parameters, roughly ten times the $\sim\!200{,}000$ of the pointwise MLP and PINN. Along the second, and the more revealing one, training time, the models split sharply. The MLP finishes in under three minutes, and the plain FNO and the two PINNs in fifteen to twenty. The PINO stands apart at well over an hour, because its physics residual is assembled on the finer $256 \times 256$ grid (Section~\ref{sec:m_pino}) at every iteration.

\begin{table}[htbp]
    \centering
    \caption{Model size, training time on the GTX~1660~SUPER, and final training loss. The final-loss column is \emph{not} comparable across rows: each objective combines data, residual and initial-condition terms differently, so its magnitude reflects the makeup of the loss rather than solution accuracy. Accuracy is measured instead by the relative-error metrics of the following sections.}
    \label{tab:training}
    \small
    \begin{tabular}{llrrr}
        \hline
        \textbf{Model} & \textbf{Regime} & \textbf{Parameters} & \textbf{Train time (s)} & \textbf{Final loss} \\
        \hline
        MLP  & data only    & $198{,}401$     & $142.8$  & $3.5 \times 10^{-5}$ \\
        FNO  & supervised   & $2{,}106{,}146$ & $1103.8$ & $3.6 \times 10^{-6}$ \\
        PINO & with data    & $2{,}106{,}146$ & $5236.2$ & $8.1 \times 10^{-5}$ \\
        PINO & no data      & $2{,}106{,}146$ & $4221.6$ & $4.7 \times 10^{-5}$ \\
        PINN & with data    & $198{,}658$     & $884.4$  & $7.9 \times 10^{-5}$ \\
        PINN & no data      & $198{,}658$     & $861.3$  & $1.2 \times 10^{-4}$ \\
        \hline
    \end{tabular}
    \\[0.4em]
    \autoriapropria
\end{table}

That hour buys generality, which is the return that justifies the operators' cost. A single FNO or PINO covers the entire $20$-point family~\eqref{eq:param_grid} from one training, whereas each MLP and PINN is welded to $\alpha = \beta = 3.21$ and would need retraining from scratch for any other value. Spread over the family, the operators, and the FNO above all, are the cheaper route to a solution at a new parameter, since inference there is a single forward pass.

One caution belongs up front, because the final-loss column invites a ranking it cannot support. The losses do not share a scale: each objective sums a different subset of data, residual and initial-condition terms. The PINO with data reaches a larger final loss ($8.1 \times 10^{-5}$) than the PINO without ($4.7 \times 10^{-5}$) because the extra supervised term adds its own residual to the total, and the sections below show it is in fact the more accurate of the two. The plain FNO carries the smallest loss in the table ($3.6 \times 10^{-6}$) because its objective is a single data term with nothing added. Training loss is a poor proxy for accuracy here, and every comparison that follows rests on the held-out relative-error metrics of Section~\ref{sec:m_metrics}.

% ===============================================================
\section{Spectral Bias, From Theory to the Network}
\label{sec:r_bias}

The models share one behaviour above all others, and before comparing them we establish its mechanism. Spectral bias, derived in Section~\ref{subsec:spectral_bias}, predicts that gradient training resolves low frequencies before high ones at a rate set by the eigenvalues of the Neural Tangent Kernel. We confirm this in two steps. First we check the prediction where it can be made exact, on the toy Fourier-feature model of Section~\ref{sec:m_probe_theory}. Then we look for the same mechanism in the kernel of a real Boussinesq PINN, where the assumptions hold only approximately.

The toy model removes every approximation, so it is the right place to test the theory quantitatively. Its Jacobian is constant by construction, so the frozen-kernel dynamics hold exactly and the closed-form predictions of Section~\ref{subsec:spectral_bias} can be compared against measurement with nothing left to interpretation. Figure~\ref{fig:theory} reports the three checks.

\begin{figure}[htbp]
    \centering
    \begin{subfigure}[t]{0.48\textwidth}
        \centering
        \includegraphics[width=\linewidth]{theory/ntk_error_decay.png}
        \caption{Per-eigendirection error decay.}
        \label{fig:theory_error}
    \end{subfigure}
    \hfill
    \begin{subfigure}[t]{0.48\textwidth}
        \centering
        \includegraphics[width=\linewidth]{theory/ntk_spectral_decay.png}
        \caption{Per-mode spectral decay.}
        \label{fig:theory_spectral}
    \end{subfigure}

    \vspace{0.8em}
    \begin{subfigure}[t]{0.60\textwidth}
        \centering
        \includegraphics[width=\linewidth]{theory/ntk_iteration_count.png}
        \caption{Measured versus predicted iteration count $n_k$.}
        \label{fig:theory_iter}
    \end{subfigure}
    \caption{Spectral-bias theory validated on the Fourier-feature model of Section~\ref{sec:m_probe_theory}, trained by full-batch gradient descent at $\mu = 1/\lambda_{\max}$ so that the linearized update~\eqref{eq:error_discrete} is exact. Measured curves (markers) overlay the theoretical predictions (lines) until the float64 floor.}
    \label{fig:theory}
    \autoriapropria
\end{figure}

Each panel confirms one prediction. In Figure~\ref{fig:theory_error} the amplitude of the error along each eigendirection follows the predicted $(1 - \mu\lambda_k)^n$ of~\eqref{eq:error_discrete}: the measured markers sit on the theoretical curves down to the float64 floor, and the directions with large eigenvalues collapse first while those with small eigenvalues linger, the per-eigendirection ordering of~\eqref{eq:error_evolution}. Figure~\ref{fig:theory_spectral} reads the same dynamics in Fourier space, where the measured per-mode amplitude tracks the prediction of~\eqref{eq:fourier_evolution} and the low wavenumbers resolve first, the frequency principle stated literally. Figure~\ref{fig:theory_iter} makes the agreement quantitative: plotting, for each mode, the measured iteration at which its amplitude first falls below the target precision against the closed-form estimate $n_k$ of~\eqref{eq:iteration_count} places the points on the $y = x$ diagonal, coloured by $\log_{10}\lambda_k$, with the smallest-eigenvalue modes demanding the most iterations. The theory holds to numerical precision on the model built to satisfy its assumptions.

With the prediction confirmed in the exact case, we look for the same mechanism in a real network, whose kernel drifts. Figure~\ref{fig:ntk_panel} shows the diagnostic of Section~\ref{sec:m_probe_ntk} on the Boussinesq PINN residual for three network widths.

\begin{figure}[htbp]
    \centering
    \begin{subfigure}[t]{0.48\textwidth}
        \centering
        \includegraphics[width=\linewidth]{pinn/ntk/pinn_ntk_panel_param.png}
        \caption{Relative parameter drift $\|\theta-\theta_0\|/\|\theta_0\|$.}
        \label{fig:ntk_param}
    \end{subfigure}
    \hfill
    \begin{subfigure}[t]{0.48\textwidth}
        \centering
        \includegraphics[width=\linewidth]{pinn/ntk/pinn_ntk_panel_kernel.png}
        \caption{Relative kernel drift $\|K-K_0\|/\|K_0\|$.}
        \label{fig:ntk_kernel}
    \end{subfigure}

    \vspace{0.8em}
    \begin{subfigure}[t]{0.48\textwidth}
        \centering
        \includegraphics[width=\linewidth]{pinn/ntk/pinn_ntk_panel_spectrum.png}
        \caption{Eigenvalue spectrum of $K$: initialization (solid) versus end of training (dashed), on log--log axes.}
        \label{fig:ntk_spectrum}
    \end{subfigure}
    \caption{Neural Tangent Kernel diagnostic for the Boussinesq PINN residual at $\alpha = \beta = 3.21$, for widths $8$, $128$ and $512$.}
    \label{fig:ntk_panel}
    \autoriapropria
\end{figure}

The diagnostic confirms both ingredients the theory depends on. The frozen-kernel assumption behind~\eqref{eq:error_linear_ode} sharpens as the network widens: the width-$512$ network drifts by only about $0.06\%$ in its parameters and about $12\%$ in its kernel over the whole run, while the width-$8$ network drifts by several percent in parameters and far more in its kernel, and the curves order themselves monotonically in width. The eigenvalue spectrum of the last panel then supplies the driver: it spans many orders of magnitude and falls off steeply, with the width-$8$ kernel collapsing to the numerical floor after roughly $45$ modes and the wider kernels decaying more gradually but still steeply. By~\eqref{eq:training_time} the time to resolve an eigendirection is inversely proportional to its eigenvalue, so this steep decay means every direction past the leading smooth ones is learned prohibitively slowly. This is the eigenvalue disparity that \textcite{wang2022ntk} name as the source of spectral bias, and it sets the expectation for everything that follows: the network should pick up the low-frequency structure readily and stall on the high-frequency detail. The exact toy confirmed the prediction; the real kernel shows the same mechanism at work, and the band-tracking curves of the next section can be read as that mechanism operating inside the full architectures.

% ===============================================================
\section{Learning by Frequency Band}
\label{sec:r_bands}

We now watch spectral bias develop as the optimizer runs. Figure~\ref{fig:spectral_bias} tracks the relative spectral error in the low, mid and high bands of~\eqref{eq:bands} across training for all six models. Within any single panel the figure tells one clean story; across panels it carries the caveat of Section~\ref{sec:m_metrics}, that the operators are monitored on the near-linear sample $\alpha = \beta = 0.1$ while the pointwise MLP and PINN are monitored on the harder $\alpha = \beta = 3.21$. Cross-family magnitudes are therefore not directly comparable in this figure, and the apples-to-apples comparison waits for Section~\ref{sec:r_spectrum}, where every model is evaluated at the common parameter $3.21$.

\begin{figure}[htbp]
    \centering
    \begin{subfigure}[t]{0.48\textwidth}
        \centering
        \includegraphics[width=\linewidth]{spectral_bias_evolution_mlp_data_only.png}
        \caption{MLP (data only), at $\alpha = \beta = 3.21$.}
        \label{fig:sb_mlp}
    \end{subfigure}
    \hfill
    \begin{subfigure}[t]{0.48\textwidth}
        \centering
        \includegraphics[width=\linewidth]{spectral_bias_evolution_pinn_no_data.png}
        \caption{PINN (no data), at $\alpha = \beta = 3.21$.}
    \end{subfigure}

    \vspace{0.6em}
    \begin{subfigure}[t]{0.48\textwidth}
        \centering
        \includegraphics[width=\linewidth]{spectral_bias_evolution_pinn_with_data.png}
        \caption{PINN (with data), at $\alpha = \beta = 3.21$.}
    \end{subfigure}
    \hfill
    \begin{subfigure}[t]{0.48\textwidth}
        \centering
        \includegraphics[width=\linewidth]{spectral_bias_evolution_fno.png}
        \caption{FNO, at $\alpha = \beta = 0.1$.}
    \end{subfigure}

    \vspace{0.6em}
    \begin{subfigure}[t]{0.48\textwidth}
        \centering
        \includegraphics[width=\linewidth]{spectral_bias_evolution_pino_no_data.png}
        \caption{PINO (no data), at $\alpha = \beta = 0.1$.}
    \end{subfigure}
    \hfill
    \begin{subfigure}[t]{0.48\textwidth}
        \centering
        \includegraphics[width=\linewidth]{spectral_bias_evolution_pino_with_data.png}
        \caption{PINO (with data), at $\alpha = \beta = 0.1$.}
    \end{subfigure}
    \caption{Relative spectral error by frequency band during training. In every panel the low band (blue) is learned first and driven lowest, the mid band (purple) trails it, and the high band (red) is learned last and least, the signature of spectral bias. Note the mixed reference samples: the pointwise MLP and PINN are tracked at $\alpha = \beta = 3.21$, the operators at $\alpha = \beta = 0.1$.}
    \label{fig:spectral_bias}
    \autoriapropria
\end{figure}

The within-panel ordering is the same in all six cases and admits no exception: the low band is driven lowest, the mid band trails, and the high band is learned last and least. Every architecture and every regime obeys this ordering, which is the measured face of the NTK spectrum of Figure~\ref{fig:ntk_panel} and the toy validation of Figure~\ref{fig:theory}. What varies from panel to panel is the size of the gap and how far each band is pushed down.

The MLP panel (Figure~\ref{fig:sb_mlp}) is the purest specimen, which is why we read it first. With no physics term and no initial-condition term competing against the data fit, its ordering is spectral bias with no confounders. The low band falls from $0.837$ at the first logged epoch to $0.015$ at the last, almost two orders of magnitude, while the mid band falls from $1.004$ to $0.114$ and the high band barely stirs, moving only from $1.009$ to $0.704$. The low frequencies are crushed and the high frequencies are left almost untouched, which is the prediction of Section~\ref{subsec:spectral_bias} in its cleanest empirical form. Every other panel is a variation on it with data or physics added.

The comparisons that \emph{are} valid across panels hold the architecture and the reference sample fixed and change only the data term, and they are decisive. For the PINN at $3.21$, adding data lowers the terminal low-band error from about $0.13$ to about $0.06$ and the mid band from about $0.55$ to about $0.18$, while leaving the high band near $0.70$ almost unchanged. For the PINO at $0.1$ the effect is larger: the data term drives the low band from about $0.37$ to about $0.02$ and the mid band from about $0.83$ to about $0.05$, again with a much smaller gain in the high band, which falls only from about $0.83$ to about $0.58$. In both families the supervised term is a strong lever on the low and mid bands and a weak one on the high band. The physics residual alone, in the two no-data panels, never brings any band near the levels reached with data, and the PINO~(no~data) high band scarcely moves over the whole run, staying above $0.83$.

The FNO panel, monitored at the same easy parameter as the PINO, reaches the lowest low- and mid-band errors in the figure (about $0.013$ and $0.016$) and the lowest high-band error too (about $0.15$). Even so, its high band sits an order of magnitude above its low band. Read across its panels, Figure~\ref{fig:spectral_bias} says that operators do not \emph{solve} spectral bias, they \emph{compress} it: the best model still leaves a high-frequency residual of roughly $15\%$, and every other model leaves more.

% ===============================================================
\section{Reproducing the Spectrum}
\label{sec:r_spectrum}

To compare the families on equal terms we fix $\alpha = \beta = 3.21$ and inspect the full spatial spectrum of the prediction at three times, $t = 0$, $t = T/2 = 15$ and $t = T = 30$. Figures~\ref{fig:spec_mlp}--\ref{fig:spec_pino_nodata} report the six models; in every panel the black curve is the reference spatial spectrum and the dashed orange curve is the prediction, with the relative-error curves and the time-averaged error at the foot of each figure.

\begin{figure}[htbp]
    \centering
    \begin{subfigure}[t]{0.48\textwidth}
        \centering
        \includegraphics[width=\linewidth]{mlp/mlp_model2_spectral_panel_spectrum_1.png}
        \caption{Spectrum, $t = 0$.}
    \end{subfigure}
    \hfill
    \begin{subfigure}[t]{0.48\textwidth}
        \centering
        \includegraphics[width=\linewidth]{mlp/mlp_model2_spectral_panel_relerr_1.png}
        \caption{Relative error, $t = 0$.}
    \end{subfigure}

    \vspace{0.6em}
    \begin{subfigure}[t]{0.48\textwidth}
        \centering
        \includegraphics[width=\linewidth]{mlp/mlp_model2_spectral_panel_spectrum_2.png}
        \caption{Spectrum, $t = 15$.}
    \end{subfigure}
    \hfill
    \begin{subfigure}[t]{0.48\textwidth}
        \centering
        \includegraphics[width=\linewidth]{mlp/mlp_model2_spectral_panel_relerr_2.png}
        \caption{Relative error, $t = 15$.}
    \end{subfigure}

    \vspace{0.6em}
    \begin{subfigure}[t]{0.48\textwidth}
        \centering
        \includegraphics[width=\linewidth]{mlp/mlp_model2_spectral_panel_spectrum_3.png}
        \caption{Spectrum, $t = 30$.}
    \end{subfigure}
    \hfill
    \begin{subfigure}[t]{0.48\textwidth}
        \centering
        \includegraphics[width=\linewidth]{mlp/mlp_model2_spectral_panel_relerr_3.png}
        \caption{Relative error, $t = 30$.}
    \end{subfigure}

    \vspace{0.6em}
    \begin{subfigure}[t]{0.68\textwidth}
        \centering
        \includegraphics[width=\linewidth]{mlp/mlp_model2_spectral_panel_mean.png}
        \caption{Mean relative spectral error over time.}
    \end{subfigure}
    \caption{MLP (data only) spatial spectrum at $\alpha = \beta = 3.21$. Fitting the reference with no equation and no physics term, the network reproduces the low-wavenumber crests but undershoots the tail: the relative error climbs with wavenumber, the purest picture of spectral bias among the six models.}
    \label{fig:spec_mlp}
    \autoriapropria
\end{figure}

\begin{figure}[htbp]
    \centering
    \begin{subfigure}[t]{0.48\textwidth}
        \centering
        \includegraphics[width=\linewidth]{pinn/no_data/pinn_no_data_model2_spectral_panel_spectrum_1.png}
        \caption{Spectrum, $t = 0$.}
    \end{subfigure}
    \hfill
    \begin{subfigure}[t]{0.48\textwidth}
        \centering
        \includegraphics[width=\linewidth]{pinn/no_data/pinn_no_data_model2_spectral_panel_relerr_1.png}
        \caption{Relative error, $t = 0$.}
    \end{subfigure}

    \vspace{0.6em}
    \begin{subfigure}[t]{0.48\textwidth}
        \centering
        \includegraphics[width=\linewidth]{pinn/no_data/pinn_no_data_model2_spectral_panel_spectrum_2.png}
        \caption{Spectrum, $t = 15$.}
    \end{subfigure}
    \hfill
    \begin{subfigure}[t]{0.48\textwidth}
        \centering
        \includegraphics[width=\linewidth]{pinn/no_data/pinn_no_data_model2_spectral_panel_relerr_2.png}
        \caption{Relative error, $t = 15$.}
    \end{subfigure}

    \vspace{0.6em}
    \begin{subfigure}[t]{0.48\textwidth}
        \centering
        \includegraphics[width=\linewidth]{pinn/no_data/pinn_no_data_model2_spectral_panel_spectrum_3.png}
        \caption{Spectrum, $t = 30$.}
    \end{subfigure}
    \hfill
    \begin{subfigure}[t]{0.48\textwidth}
        \centering
        \includegraphics[width=\linewidth]{pinn/no_data/pinn_no_data_model2_spectral_panel_relerr_3.png}
        \caption{Relative error, $t = 30$.}
    \end{subfigure}

    \vspace{0.6em}
    \begin{subfigure}[t]{0.68\textwidth}
        \centering
        \includegraphics[width=\linewidth]{pinn/no_data/pinn_no_data_model2_spectral_panel_mean.png}
        \caption{Mean relative spectral error over time.}
    \end{subfigure}
    \caption{PINN (no data) spatial spectrum at $\alpha = \beta = 3.21$. The network cannot represent the $\sech^2$ tail even at $t = 0$, and the relative error rises monotonically with wavenumber; the mean spectral error grows from $0.17$ at $t = 0$ to $0.70$ at $t = 30$.}
    \label{fig:spec_pinn_nodata}
    \autoriapropria
\end{figure}

\begin{figure}[htbp]
    \centering
    \begin{subfigure}[t]{0.48\textwidth}
        \centering
        \includegraphics[width=\linewidth]{pinn/with_data/pinn_model2_spectral_panel_spectrum_1.png}
        \caption{Spectrum, $t = 0$.}
    \end{subfigure}
    \hfill
    \begin{subfigure}[t]{0.48\textwidth}
        \centering
        \includegraphics[width=\linewidth]{pinn/with_data/pinn_model2_spectral_panel_relerr_1.png}
        \caption{Relative error, $t = 0$.}
    \end{subfigure}

    \vspace{0.6em}
    \begin{subfigure}[t]{0.48\textwidth}
        \centering
        \includegraphics[width=\linewidth]{pinn/with_data/pinn_model2_spectral_panel_spectrum_2.png}
        \caption{Spectrum, $t = 15$.}
    \end{subfigure}
    \hfill
    \begin{subfigure}[t]{0.48\textwidth}
        \centering
        \includegraphics[width=\linewidth]{pinn/with_data/pinn_model2_spectral_panel_relerr_2.png}
        \caption{Relative error, $t = 15$.}
    \end{subfigure}

    \vspace{0.6em}
    \begin{subfigure}[t]{0.48\textwidth}
        \centering
        \includegraphics[width=\linewidth]{pinn/with_data/pinn_model2_spectral_panel_spectrum_3.png}
        \caption{Spectrum, $t = 30$.}
    \end{subfigure}
    \hfill
    \begin{subfigure}[t]{0.48\textwidth}
        \centering
        \includegraphics[width=\linewidth]{pinn/with_data/pinn_model2_spectral_panel_relerr_3.png}
        \caption{Relative error, $t = 30$.}
    \end{subfigure}

    \vspace{0.6em}
    \begin{subfigure}[t]{0.68\textwidth}
        \centering
        \includegraphics[width=\linewidth]{pinn/with_data/pinn_model2_spectral_panel_mean.png}
        \caption{Mean relative spectral error over time.}
    \end{subfigure}
    \caption{PINN (with data) spatial spectrum at $\alpha = \beta = 3.21$. The low-wavenumber crests are tracked more tightly than in Figure~\ref{fig:spec_pinn_nodata}, but the high-wavenumber undershoot and the growth of the mean error in time ($0.17 \to 0.58 \to 0.68$) remain.}
    \label{fig:spec_pinn_data}
    \autoriapropria
\end{figure}

The pointwise panels lay the pathology bare, and they carry it into the initial state. Both PINN regimes already fail to represent the high-wavenumber content of the $\sech^2$ profile at $t = 0$, and the relative error climbs toward unity as the mode index grows. Without data (Figure~\ref{fig:spec_pinn_nodata}) the predicted spectrum has shed most of its mid- and high-frequency amplitude by $t = 30$, and the mean spectral error reaches $0.70$. Adding data (Figure~\ref{fig:spec_pinn_data}) tightens the low-wavenumber crests, in step with the low-band gain of Figure~\ref{fig:spectral_bias}, yet the aggregate mean barely moves ($0.17 \to 0.58 \to 0.68$), because the mid and high modes dominate it and the data term fails to recover them at this hard parameter. The MLP (Figure~\ref{fig:spec_mlp}) tells the same story with the physics stripped away: it fits the reference where the frequencies are low and undershoots where they are high, its mean spectral error rising from about $0.16$ at $t = 0$ to about $0.51$ at $t = 15$ and about $0.65$ at $t = 30$. Data helps the pointwise models where the kernel is favourable and does not rescue them where it is adverse.

\begin{figure}[htbp]
    \centering
    \begin{subfigure}[t]{0.48\textwidth}
        \centering
        \includegraphics[width=\linewidth]{fno/fno_model2_spectral_panel_spectrum_1.png}
        \caption{Spectrum, $t = 0$.}
    \end{subfigure}
    \hfill
    \begin{subfigure}[t]{0.48\textwidth}
        \centering
        \includegraphics[width=\linewidth]{fno/fno_model2_spectral_panel_relerr_1.png}
        \caption{Relative error, $t = 0$.}
    \end{subfigure}

    \vspace{0.6em}
    \begin{subfigure}[t]{0.48\textwidth}
        \centering
        \includegraphics[width=\linewidth]{fno/fno_model2_spectral_panel_spectrum_2.png}
        \caption{Spectrum, $t = 15$.}
    \end{subfigure}
    \hfill
    \begin{subfigure}[t]{0.48\textwidth}
        \centering
        \includegraphics[width=\linewidth]{fno/fno_model2_spectral_panel_relerr_2.png}
        \caption{Relative error, $t = 15$.}
    \end{subfigure}

    \vspace{0.6em}
    \begin{subfigure}[t]{0.48\textwidth}
        \centering
        \includegraphics[width=\linewidth]{fno/fno_model2_spectral_panel_spectrum_3.png}
        \caption{Spectrum, $t = 30$.}
    \end{subfigure}
    \hfill
    \begin{subfigure}[t]{0.48\textwidth}
        \centering
        \includegraphics[width=\linewidth]{fno/fno_model2_spectral_panel_relerr_3.png}
        \caption{Relative error, $t = 30$.}
    \end{subfigure}

    \vspace{0.6em}
    \begin{subfigure}[t]{0.68\textwidth}
        \centering
        \includegraphics[width=\linewidth]{fno/fno_model2_spectral_panel_mean.png}
        \caption{Mean relative spectral error over time.}
    \end{subfigure}
    \caption{FNO spatial spectrum at $\alpha = \beta = 3.21$. Low and mid wavenumbers are tracked almost exactly (mean error $0.013$ at $t = 0$), but the error migrates to the high modes and grows steadily in time, reaching a mean of $0.40$ at $t = 30$.}
    \label{fig:spec_fno}
    \autoriapropria
\end{figure}

\begin{figure}[htbp]
    \centering
    \begin{subfigure}[t]{0.48\textwidth}
        \centering
        \includegraphics[width=\linewidth]{pino/with_data/pino_model2_spectral_panel_spectrum_1.png}
        \caption{Spectrum, $t = 0$.}
    \end{subfigure}
    \hfill
    \begin{subfigure}[t]{0.48\textwidth}
        \centering
        \includegraphics[width=\linewidth]{pino/with_data/pino_model2_spectral_panel_relerr_1.png}
        \caption{Relative error, $t = 0$.}
    \end{subfigure}

    \vspace{0.6em}
    \begin{subfigure}[t]{0.48\textwidth}
        \centering
        \includegraphics[width=\linewidth]{pino/with_data/pino_model2_spectral_panel_spectrum_2.png}
        \caption{Spectrum, $t = 15$.}
    \end{subfigure}
    \hfill
    \begin{subfigure}[t]{0.48\textwidth}
        \centering
        \includegraphics[width=\linewidth]{pino/with_data/pino_model2_spectral_panel_relerr_2.png}
        \caption{Relative error, $t = 15$.}
    \end{subfigure}

    \vspace{0.6em}
    \begin{subfigure}[t]{0.48\textwidth}
        \centering
        \includegraphics[width=\linewidth]{pino/with_data/pino_model2_spectral_panel_spectrum_3.png}
        \caption{Spectrum, $t = 30$.}
    \end{subfigure}
    \hfill
    \begin{subfigure}[t]{0.48\textwidth}
        \centering
        \includegraphics[width=\linewidth]{pino/with_data/pino_model2_spectral_panel_relerr_3.png}
        \caption{Relative error, $t = 30$.}
    \end{subfigure}

    \vspace{0.6em}
    \begin{subfigure}[t]{0.68\textwidth}
        \centering
        \includegraphics[width=\linewidth]{pino/with_data/pino_model2_spectral_panel_mean.png}
        \caption{Mean relative spectral error over time.}
    \end{subfigure}
    \caption{PINO (with data) spatial spectrum at $\alpha = \beta = 3.21$. The prediction tracks the reference across the band and, unlike the FNO, the mean error stays bounded in time ($0.077 \to 0.13 \to 0.20$): the physics residual steadies the temporal evolution.}
    \label{fig:spec_pino_data}
    \autoriapropria
\end{figure}

\begin{figure}[htbp]
    \centering
    \begin{subfigure}[t]{0.48\textwidth}
        \centering
        \includegraphics[width=\linewidth]{pino/no_data/pino_no_data_model2_spectral_panel_spectrum_1.png}
        \caption{Spectrum, $t = 0$.}
    \end{subfigure}
    \hfill
    \begin{subfigure}[t]{0.48\textwidth}
        \centering
        \includegraphics[width=\linewidth]{pino/no_data/pino_no_data_model2_spectral_panel_relerr_1.png}
        \caption{Relative error, $t = 0$.}
    \end{subfigure}

    \vspace{0.6em}
    \begin{subfigure}[t]{0.48\textwidth}
        \centering
        \includegraphics[width=\linewidth]{pino/no_data/pino_no_data_model2_spectral_panel_spectrum_2.png}
        \caption{Spectrum, $t = 15$.}
    \end{subfigure}
    \hfill
    \begin{subfigure}[t]{0.48\textwidth}
        \centering
        \includegraphics[width=\linewidth]{pino/no_data/pino_no_data_model2_spectral_panel_relerr_2.png}
        \caption{Relative error, $t = 15$.}
    \end{subfigure}

    \vspace{0.6em}
    \begin{subfigure}[t]{0.48\textwidth}
        \centering
        \includegraphics[width=\linewidth]{pino/no_data/pino_no_data_model2_spectral_panel_spectrum_3.png}
        \caption{Spectrum, $t = 30$.}
    \end{subfigure}
    \hfill
    \begin{subfigure}[t]{0.48\textwidth}
        \centering
        \includegraphics[width=\linewidth]{pino/no_data/pino_no_data_model2_spectral_panel_relerr_3.png}
        \caption{Relative error, $t = 30$.}
    \end{subfigure}

    \vspace{0.6em}
    \begin{subfigure}[t]{0.68\textwidth}
        \centering
        \includegraphics[width=\linewidth]{pino/no_data/pino_no_data_model2_spectral_panel_mean.png}
        \caption{Mean relative spectral error over time.}
    \end{subfigure}
    \caption{PINO (no data) spatial spectrum at $\alpha = \beta = 3.21$. Without the supervised term the low- and mid-band undershoot returns (mean error $0.11 \to 0.39 \to 0.33$), confirming that the physics residual alone does not overcome spectral bias.}
    \label{fig:spec_pino_nodata}
    \autoriapropria
\end{figure}

The operators, learning in the frequency domain, invert this picture at low and mid wavenumbers. The FNO (Figure~\ref{fig:spec_fno}) reproduces the low and mid spectrum almost exactly at early times, a mean error of $0.013$ at $t = 0$, an order of magnitude below the pointwise models, because those modes are precisely what its spectral layers parametrize. Its weakness lies in time: the error gathers in the high modes and \emph{grows}, from a mean of $0.013$ at $t = 0$ to $0.40$ at $t = 30$. The PINO with data (Figure~\ref{fig:spec_pino_data}) inherits the operator's spectral accuracy without this temporal decay, its mean error rising only from $0.077$ to $0.20$ over the same horizon, the steadiest of the six at late times. The contrast is sharp at $t = 30$, where the PINO with data ($0.20$) clearly beats the FNO ($0.40$) though the FNO was ahead at $t = 0$: the physics residual trades a slightly worse start for a much steadier trajectory. The PINO without data (Figure~\ref{fig:spec_pino_nodata}) forfeits the advantage, its low- and mid-band undershoot returning and its mean error climbing to $0.33$--$0.39$, midway between the operators-with-data and the PINNs. The physics constraint by itself does not replace data.

% ===============================================================
\section{Where Physics Earns Its Keep}
\label{sec:r_gibbs}

The temporal growth of the FNO error in Figure~\ref{fig:spec_fno} has a specific and instructive cause, and isolating it shows the one place where the physics term pays for itself. As Section~\ref{sec:m_fno} explains, the operator applies the Fourier transform along the temporal axis as well as the spatial one, implicitly assuming the trajectory is periodic in time. It is not: the field at $t = T$ does not match the field at $t = 0$, so the assumption manufactures an artificial discontinuity at the temporal boundary and excites the Gibbs phenomenon there \parencite{gibbs1899}. Figures~\ref{fig:ab_fno} and~\ref{fig:ab_pino} make the effect and its cure visible.

\begin{figure}[htbp]
    \centering
    \begin{subfigure}[t]{0.48\textwidth}
        \centering
        \includegraphics[width=\linewidth]{fno/fno_model2_alpha_beta_panel_surface_1.png}
        \caption{Predicted $\eta(x,t)$, $\alpha=\beta=0.1$.}
    \end{subfigure}
    \hfill
    \begin{subfigure}[t]{0.48\textwidth}
        \centering
        \includegraphics[width=\linewidth]{fno/fno_model2_alpha_beta_panel_relerr_1.png}
        \caption{Relative error, $\alpha=\beta=0.1$.}
    \end{subfigure}

    \vspace{0.6em}
    \begin{subfigure}[t]{0.48\textwidth}
        \centering
        \includegraphics[width=\linewidth]{fno/fno_model2_alpha_beta_panel_surface_2.png}
        \caption{Predicted $\eta(x,t)$, $\alpha=\beta=3.21$.}
    \end{subfigure}
    \hfill
    \begin{subfigure}[t]{0.48\textwidth}
        \centering
        \includegraphics[width=\linewidth]{fno/fno_model2_alpha_beta_panel_relerr_2.png}
        \caption{Relative error, $\alpha=\beta=3.21$.}
    \end{subfigure}

    \vspace{0.6em}
    \begin{subfigure}[t]{0.48\textwidth}
        \centering
        \includegraphics[width=\linewidth]{fno/fno_model2_alpha_beta_panel_surface_3.png}
        \caption{Predicted $\eta(x,t)$, $\alpha=\beta=4.2$.}
    \end{subfigure}
    \hfill
    \begin{subfigure}[t]{0.48\textwidth}
        \centering
        \includegraphics[width=\linewidth]{fno/fno_model2_alpha_beta_panel_relerr_3.png}
        \caption{Relative error, $\alpha=\beta=4.2$.}
    \end{subfigure}

    \vspace{0.6em}
    \begin{subfigure}[t]{0.68\textwidth}
        \centering
        \includegraphics[width=\linewidth]{fno/fno_model2_alpha_beta_panel_box.png}
        \caption{Relative error distribution vs.\ $\alpha=\beta$.}
    \end{subfigure}
    \caption{FNO across the nonlinearity axis ($\alpha=\beta \in \{0.1, 3.21, 4.2\}$), evaluated at resolution $256$. The predicted surfaces (left) capture the two counter-propagating pulses, but the time-resolved relative error (right) spikes sharply at the final time $t = 30$: the temporal Gibbs artifact of treating time as a periodic axis.}
    \label{fig:ab_fno}
    \autoriapropria
\end{figure}

\begin{figure}[htbp]
    \centering
    \begin{subfigure}[t]{0.48\textwidth}
        \centering
        \includegraphics[width=\linewidth]{pino/with_data/pino_model2_alpha_beta_panel_surface_1.png}
        \caption{Predicted $\eta(x,t)$, $\alpha=\beta=0.1$.}
    \end{subfigure}
    \hfill
    \begin{subfigure}[t]{0.48\textwidth}
        \centering
        \includegraphics[width=\linewidth]{pino/with_data/pino_model2_alpha_beta_panel_relerr_1.png}
        \caption{Relative error, $\alpha=\beta=0.1$.}
    \end{subfigure}

    \vspace{0.6em}
    \begin{subfigure}[t]{0.48\textwidth}
        \centering
        \includegraphics[width=\linewidth]{pino/with_data/pino_model2_alpha_beta_panel_surface_2.png}
        \caption{Predicted $\eta(x,t)$, $\alpha=\beta=3.21$.}
    \end{subfigure}
    \hfill
    \begin{subfigure}[t]{0.48\textwidth}
        \centering
        \includegraphics[width=\linewidth]{pino/with_data/pino_model2_alpha_beta_panel_relerr_2.png}
        \caption{Relative error, $\alpha=\beta=3.21$.}
    \end{subfigure}

    \vspace{0.6em}
    \begin{subfigure}[t]{0.48\textwidth}
        \centering
        \includegraphics[width=\linewidth]{pino/with_data/pino_model2_alpha_beta_panel_surface_3.png}
        \caption{Predicted $\eta(x,t)$, $\alpha=\beta=4.2$.}
    \end{subfigure}
    \hfill
    \begin{subfigure}[t]{0.48\textwidth}
        \centering
        \includegraphics[width=\linewidth]{pino/with_data/pino_model2_alpha_beta_panel_relerr_3.png}
        \caption{Relative error, $\alpha=\beta=4.2$.}
    \end{subfigure}

    \vspace{0.6em}
    \begin{subfigure}[t]{0.68\textwidth}
        \centering
        \includegraphics[width=\linewidth]{pino/with_data/pino_model2_alpha_beta_panel_box.png}
        \caption{Relative error distribution vs.\ $\alpha=\beta$.}
    \end{subfigure}
    \caption{PINO (with data) across the same nonlinearity axis and resolution. The predicted surfaces are visually indistinguishable from the FNO, but the terminal-time spike of Figure~\ref{fig:ab_fno} is gone: the error rises only gently toward $t = 30$. The physics residual and the initial-condition term act as a temporal regularizer.}
    \label{fig:ab_pino}
    \autoriapropria
\end{figure}

In the FNO panel (Figure~\ref{fig:ab_fno}) every parameter shows the same feature: after an initial transient the time-resolved error settles to a plateau and then jumps at $t = 30$ to two or three times that plateau. The spike appears at all three nonlinearities and is the dominant part of the FNO's late-time error. In the PINO panel (Figure~\ref{fig:ab_pino}) the predicted surfaces are visually identical, yet the spike is gone and the error rises only slightly toward the boundary. The PDE residual and the initial-condition term penalize exactly the temporal inconsistency that a periodic transform would otherwise tolerate, damping the boundary oscillation. This is the one place in the study where the physics term delivers a clear, concrete benefit: it suppresses the FNO's Gibbs artifact. The benefit is a suppression and not an elimination, since a small terminal rise remains, and it arrives at the modelling cost of the residual machinery.

% ===============================================================
\section{Two Stress Tests}
\label{sec:r_stress}

The comparisons so far held the operators on their training conditions. Two stress tests now push them off those conditions, one along the nonlinearity axis and one along the query resolution. Both concern the operators alone, since each pointwise MLP and PINN is trained at a single parameter and a single resolution and has nothing to say about either sweep.

The first test sweeps the nonlinearity. The box plots at the foot of Figures~\ref{fig:ab_fno} and~\ref{fig:ab_pino} summarize the distribution of the time-resolved error~\eqref{eq:rel_time} over the whole trajectory at each of the three parameters. Their most striking feature runs against intuition: the \emph{smallest} nonlinearity, $\alpha = \beta = 0.1$, carries the \emph{largest} relative error for both operators, a median near $0.18$ against roughly $0.12$--$0.13$ at the more nonlinear values $3.21$ and $4.2$. We read this carefully. Two effects plausibly combine. First, $\alpha = \beta = 0.1$ sits at the extreme lower edge of the training grid~\eqref{eq:param_grid}, where the operator has the least surrounding data to interpolate from, so a boundary-of-range penalty is expected. Second, the relative metric~\eqref{eq:rel_time} divides by the instantaneous norm of the reference: in the near-linear regime the solution is dominated by the two clean, slowly deforming pulses of the wave-equation limit, whose energy is concentrated and whose norm is modest, so a fixed absolute error reads as a larger relative one. The nonlinear cases spread energy across the domain and are, in relative terms, more forgiving. The box plots do not show the operators growing uniformly more accurate, in absolute terms, as nonlinearity rises. They show only that the relative error, under this normalization and this training grid, is non-monotone in the parameter, with the near-linear boundary case the least well served.

The second test queries the operator off its training resolution, probing the discretization-invariance operators are meant to enjoy. Figure~\ref{fig:res_pino} takes the PINO with data, trained at resolution $128$, and evaluates it at $128$, $256$ and $512$ at the fixed parameter $\alpha = \beta = 3.21$ with no retraining.

\begin{figure}[htbp]
    \centering
    \begin{subfigure}[t]{0.48\textwidth}
        \centering
        \includegraphics[width=\linewidth]{pino/with_data/pino_model2_resolution_panel_surface_1.png}
        \caption{Predicted $\eta(x,t)$, res.\ $128$.}
    \end{subfigure}
    \hfill
    \begin{subfigure}[t]{0.48\textwidth}
        \centering
        \includegraphics[width=\linewidth]{pino/with_data/pino_model2_resolution_panel_relerr_1.png}
        \caption{Relative error, res.\ $128$.}
    \end{subfigure}

    \vspace{0.6em}
    \begin{subfigure}[t]{0.48\textwidth}
        \centering
        \includegraphics[width=\linewidth]{pino/with_data/pino_model2_resolution_panel_surface_2.png}
        \caption{Predicted $\eta(x,t)$, res.\ $256$.}
    \end{subfigure}
    \hfill
    \begin{subfigure}[t]{0.48\textwidth}
        \centering
        \includegraphics[width=\linewidth]{pino/with_data/pino_model2_resolution_panel_relerr_2.png}
        \caption{Relative error, res.\ $256$.}
    \end{subfigure}

    \vspace{0.6em}
    \begin{subfigure}[t]{0.48\textwidth}
        \centering
        \includegraphics[width=\linewidth]{pino/with_data/pino_model2_resolution_panel_surface_3.png}
        \caption{Predicted $\eta(x,t)$, res.\ $512$.}
    \end{subfigure}
    \hfill
    \begin{subfigure}[t]{0.48\textwidth}
        \centering
        \includegraphics[width=\linewidth]{pino/with_data/pino_model2_resolution_panel_relerr_3.png}
        \caption{Relative error, res.\ $512$.}
    \end{subfigure}

    \vspace{0.6em}
    \begin{subfigure}[t]{0.68\textwidth}
        \centering
        \includegraphics[width=\linewidth]{pino/with_data/pino_model2_resolution_panel_box.png}
        \caption{Relative error distribution vs.\ resolution.}
    \end{subfigure}
    \caption{PINO (with data) queried zero-shot at resolutions $128$, $256$ and $512$ at $\alpha = \beta = 3.21$. The operator returns a coherent field at every resolution, but the relative error against the pseudospectral reference \emph{grows} with resolution (median rising from about $0.07$ at $128$ to about $0.16$ at $512$): invariance of form is not invariance of accuracy.}
    \label{fig:res_pino}
    \autoriapropria
\end{figure}

The geometric promise holds. At every resolution the operator returns a smooth, coherent surface with the correct pulse structure, with no retraining and no architectural change; the input is simply resampled on the finer grid. The accuracy promise is subtler. The box plot shows the relative error \emph{increasing} with resolution, from a median near $0.07$ at the training resolution of $128$ to about $0.12$ at $256$ and $0.16$ at $512$. The cause is that the finer references contain genuine high-wavenumber content that the operator, trained at $128$ and keeping only $16$ Fourier modes per axis, never learned to produce, so querying it against a sharper reference exposes the fine scales it is missing. A secondary effect appears at the training resolution itself, where the error oscillates in time in step with the wave reflections, the aliasing signature of a coarse grid. Zero-shot superresolution is real as a \emph{geometric} property, since the operator accepts and returns fields at any resolution, and it is not an \emph{accuracy} property: fidelity is anchored to the spectral band seen in training, and a richer grid can only reveal, never recover, the modes left out.

% ===============================================================
\section{What the Experiments Show}
\label{sec:r_synthesis}

Taken together, the experiments do not crown a single winner; they chart where each method succeeds and where it fails.

\begin{itemize}
    \item \textbf{Spectral bias is universal, and provably so here.} The toy validation of Section~\ref{sec:r_bias} matches the predicted per-mode decay and iteration counts to numerical precision, the real PINN kernel shows the same eigenvalue disparity, and every architecture then obeys the same ordering: low wavenumbers first and best, high wavenumbers last and least (Figures~\ref{fig:ntk_panel}, \ref{fig:theory}, \ref{fig:spectral_bias}). The data-only MLP is the purest instance, driving its low band by nearly two orders of magnitude while its high band barely moves. No method drives the high band below roughly $0.15$ relative error, and the best there is the plain FNO.

    \item \textbf{Data is the decisive lever, and physics is not a substitute for it.} Within both the PINN and the PINO families, adding the supervised term sharply cuts the low- and mid-band error, while the physics-only regimes revert toward the biased baseline (Figures~\ref{fig:spectral_bias}, \ref{fig:spec_pino_data}, \ref{fig:spec_pino_nodata}). Physics alone neither matches data nor removes the bias.

    \item \textbf{Operators dominate at low and mid wavenumbers.} At the common parameter $3.21$ the FNO and the PINO with data track the low and mid spectrum an order of magnitude more accurately than either PINN or the MLP (Section~\ref{sec:r_spectrum}), and they cover the whole family from one training, whereas each pointwise model is bound to a single parameter.

    \item \textbf{Physics buys temporal stability.} The FNO's error grows in time and spikes at the temporal boundary through the Gibbs artifact of its periodic time axis, and the PINO's residual suppresses the spike and keeps the late-time error bounded (Section~\ref{sec:r_gibbs}). This is the clearest concrete benefit of the physics term in this study.

    \item \textbf{The remaining limitations are real.} The near-linear boundary parameter is the least accurate (Section~\ref{sec:r_stress}), zero-shot superresolution degrades in accuracy as the query grid sharpens (Section~\ref{sec:r_stress}), and the high-frequency band is never fully resolved by any method. These are consistent features of the design, not artefacts of a single run.
\end{itemize}

The picture matches the one anticipated in Chapter~\ref{chap:intro}: these methods complement the pseudospectral solver that sets the accuracy ceiling here rather than replacing it. Their value is in amortizing the cost of the parameter sweep and in the temporal regularization that physics supplies, and it is least of all at the high wavenumbers where the reference solver stays unchallenged.
